\begin{center}
    \huge Nilpotent Groups Cheat Sheet
\end{center}

\subsection*{Commutators}

Given elements of a group $x,y \in G$, their \emph{commutator} is $[x,y] \coloneqq xyx\inv y\inv$. Given two subgroups $H_1, H_2 < G$ (or sometimes even just subsets), their commutator is the subgroup generated by all $[h_1,h_2]$ for $h_1 \in H_1$, $h_2 \in H_2$. For example, let $G = S_4$, $H_1 = \genset{(12)}$, and $H_2 = \genset{(234)}$. In this case there are only three possible commutators of the form $h_1 h_2 h_1\inv h_2\inv$: $e$, $(132)$, and $(142)$. One can check that these two elements generate $A_4$, so $[H_1,H_2] = A_4$. (Note that the \emph{set} of commutators of the form $h_1 h_2 h_1\inv h_2\inv$ is not itself a group in this case, which is why we have to say ``subgroup generated by''.)

The subgroup $[G,G]$ is called the \emph{commutator subgroup} or the \emph{derived subgroup}. It is always normal, and has the property that $G/[G,G]$ is always abelian, and if $N$ is another normal subgroup so that $G/N$ is abelian, then $N \supset [G,G]$. The group $G/[G,G]$ is called the \emph{abelianization} of $G$, written $G\ab$.

If you know $G$ is generated by a set $S$, you can find generators for $[G,G]$ in terms of commutators coming from $S$. It's not quite enough to just take the commutators of the form $s_1 s_2 s_1\inv s_2\inv$, you also have to take all conjugates of those commutators. (Idk if there's a similar statement for subgroups of the form $[H_1,H_2]$.)

\subsection*{Upper and Lower Central Series}
% The example given in this section is the modular maximal-cyclic group M_4(2), obtained as the semidirect product of C_2 acting on C_8 via multiplication by 3. This is the smallest group for which the upper and lower central series are not equal. These facts and the representation below were found at https://people.maths.bris.ac.uk/~matyd/GroupNames/1/M4(2).html

The \emph{lower central series} of $G$ is the descending series of subgroups $G = G_0 \triangleright G_1 \triangleright G_2 \triangleright \dots$, where $G_{i+1} = [G,G_i]$. For example, let $G$ be the subgroup of $\GL_2(\F_5)$ generated by the two matrices
\[\alpha = \twomatrix{0}{2}{1}{0} \qquad \text{and} \qquad \beta = \twomatrix{-1}{0}{0}{1}.\]
We will compute the lower central series of $G$. Since $\alpha$ and $\beta$ generate $G$, to compute $[G,G]$ we only need to compute all the conjugates of $\alpha \beta \alpha\inv \beta\inv$. A little calculation shows that $\alpha \beta \alpha\inv \beta\inv = -I$. Since $-I \in Z(G)$, its only conjugate is itself, so $G_1 = [G,G] = \genset{-I} \cong C_2$ a cyclic group of order 2. Also, since $-I \in Z(G)$, $G_2 = [G,G_1] = \set{I}$ the trivial group. Thus, for this group, the lower central series is $G \triangleright \genset{-I} \triangleright \set{I}$.

The \emph{upper central series} of $G$ is the ascending series of subgroups $\set{e} = Z_0 \triangleleft Z_1 \triangleleft Z_2 \triangleleft \dots$, where $Z_{i+1}$ is defined so that $Z_{i+1}/Z_i = Z(G/Z_i)$. Let's compute the upper central series of our group $G$ from above. We start by computing $Z_1 = Z(G)$. Note that $\alpha^2 = 2I \in Z(G)$, and $2I$ generates all of the diagonal matrices mod 5. We can show that any element of $Z(G)$ must be diagonal in a typical way for groups of matrices. Any matrix $\twomatrix{a}{b}{c}{d}$ which commutes with $\alpha$ must have $a = d$ and $b = 2c$. On the other hand, any matrix which commutes with $\beta$ must have $b = -b$, so $b = 0$. Thus, any matrix commuting with both $\alpha$ and $\beta$ has $a = d$ and $b = c = 0$, which is to say, it is diagonal. Hence, $Z(G)$ consists of the diagonal matrices, which is a cyclic subgroup of order 4. One can check that $G$ has order 16, so $G/Z(G)$ has order 4, and hence $G/Z(G)$ is abelian, so $Z_2 = G$. Thus, for this group, the upper central series is $\set{I} \triangleleft \genset{2I} \triangleleft G$.

Notice that in the examples, the upper and lower central series had the same length, and that every term of the upper central series contained the corresponding group in the lower central series.

\subsection*{Nilpotent Groups}

A group is \emph{nilpotent} if either of the equivalent conditions occurs:
\begin{enumerate}[(i)]
    \item the lower central series terminates in the trivial group $\set{e}$; or,
    \item the upper central series terminates in the whole group $G$.
\end{enumerate}
The \emph{nilpotency class} of $G$ is the length (number of $\triangleright$s) of either series. A group is nilpotent of class 1 iff it is abelian, so nilpotent groups are a kind of generalization of abelian groups. Our calculations above show that the group we looked at is nilpotent of class 2. Any $p$-group (group whose order is a prime power) is automatically nilpotent, and it turns out another equivalent condition for being nilpotent is:
\begin{enumerate}[(i),resume]
    \item $G$ is a direct produt of $p$-groups.
\end{enumerate}

Some other examples of nilpotent groups include the dihedral groups $D_{2^n}$, and the quaternion group $Q$ (because they are both 2-groups). Another family of examples are the \emph{Heisenberg groups}, which consist of $3 \times 3$ upper triangular matrices over $\F_q$ with 1s on the diagonal.

An example of a non-nilpotent group is $S_n$ for $n \geq 3$, because $S_n$ is centerless, so the upper central series never even gets off the ground. On the other hand, one can check that $[S_n,S_n] = A_n$, and $[S_n,A_n] = A_n$, so the lower central series terminates in $A_n$ after just one step. Notice that in the case of a non-nilpotent group, the lengths of the upper and lower central series can be different.